\section{Proofs for Section~\ref{sec:boundary}}\label{app:boundary}

Write $e=e_\tau(H)=\sgm(H/\tau)=\Prob(A=1\mid H)$ and
$\omega=e(1-e)=\Omega(H/\tau)$ with $\Omega(u)=\exp(u)/\{1+\exp(u)\}^2$, and
$m_\tau=\E\omega$. We use
\begin{gather*}
\int \Omega=1,\qquad \int|u|\Omega(u)\,du=2\log2,\qquad \int \Omega(u)\sgm(u)\,du=\tfrac12,\\
\Omega(u)\le \exp(-|u|),\qquad \int_0^\infty\sgm(-v)\,dv=\log2,\qquad \int_0^\infty v\,\sgm(-v)\,dv=\tfrac{\pi^2}{12}.
\end{gather*}

\subsection{Proof of Theorem~\ref{thm:boundary}}

\paragraph{Step 1 (kernel integrals).} For $\varrho$ bounded and continuous at $0$,
$j\ge1$ and $r,s\ge0$, substituting $h=\tau u$ and using dominated convergence
with dominating function $\|\varrho\|_\infty\bar f\Omega(u)$,
\[
\E[\Omega(H/\tau)^j(1-e)^r\,e^s\,\varrho(H)]=\tau f(0)\varrho(0)\int \Omega^j\sgm(-u)^r\sgm(u)^s\,du+o(\tau).
\]
In particular $m_\tau=\tau f(0)(1+o(1))$.

\paragraph{Step 2 (bias).} $\beta_{\ov,\tau}-\beta_0=\E[\omega\{c(H)-c(0)\}]/m_\tau$.
On $|H|\le u_0$, $|c(h)-c(0)|\le2L|h|$ and
$\E[\omega|H|]\le\bar f\tau^2\int|u|\Omega=2\log2\,\bar f\tau^2$. On $|H|>u_0$,
$\omega\le \exp(-u_0/\tau)$ and $|c(h)-c(0)|\le4B$. Hence
$|\beta_{\ov,\tau}-\beta_0|\le(4\log2\,L\bar f\tau^2+4B\exp(-u_0/\tau))/m_\tau
=(4\log2\,L\bar f/f(0))\tau(1+o(1))$.

\paragraph{Step 3 (linearization).} Since $\E[w_1\mid H]=\omega$ and
$\E[w_1Y\mid H]=\omega\mu_1$, $\E[w_1(Y-\bar\mu_1)]=0$, and likewise for arm
$0$. Define $T_a=n^{-1}\sum_iw_{ai}(Y_i-\bar\mu_a)/m_\tau$, so that
$\hat m_a-\bar\mu_a=(m_\tau/\hat D_a)T_a$, $T_1-T_0=n^{-1}\sum_i\xi_i$ and
$\E\xi=0$. Because $w_1w_0=0$ and $\E[A(1-e)^2\mid H]=\omega(1-e)$,
\begin{equation}\label{eq:sn}
s_n^2=m_\tau^{-2}\E\big[\omega\{(1-e)(\sigma_1^2+(\mu_1-\bar\mu_1)^2)+e(\sigma_0^2+(\mu_0-\bar\mu_0)^2)\}\big].
\end{equation}
By Step 1, $\bar\mu_a\to\mu_a(0)$; the terms in $(\mu_a-\bar\mu_a)^2$ are
$o(\tau)$ by dominated convergence; and the variance terms equal
$\tau f(0)\sigma_a^2(0)/2+o(\tau)$. Hence
$s_n^2=\{\sigma_1^2(0)+\sigma_0^2(0)\}/\{2\tau f(0)\}(1+o(1))$.

\paragraph{Step 4 (central limit theorem).} Since $0\le w_a\le1$, $w_a^4\le
w_a$, and $\E[(Y-\bar\mu_a)^4\mid A,H]\le C(M_4,B)$, we get
$\E[w_a^4(Y-\bar\mu_a)^4]\le Cm_\tau$ and $\E\xi^4\le C'm_\tau^{-3}$. The
Lyapunov ratio $n^{-1}\E\xi^4/s_n^4=O((n\tau)^{-1})\to0$, so
$n^{-1/2}\sum\xi_i/s_n\Rightarrow N(0,1)$. Also
$\Var(\hat D_a)/m_\tau^2\le1/(nm_\tau)\to0$, so $\hat D_a/m_\tau\to1$ in
probability and $m_\tau/\hat D_a-1=O_p((n\tau)^{-1/2})$. Now
\begin{equation}\label{eq:decomp}
\hat\beta-\beta_{\ov,\tau}=(T_1-T_0)+(m_\tau/\hat D_1-1)T_1-(m_\tau/\hat D_0-1)T_0,
\end{equation}
with $T_a=O_p((n\tau)^{-1/2})$. The remainder is
$O_p((n\tau)^{-1})=o_p(s_n/\sqrt n)$, and Slutsky's lemma gives (i). The event
that a weight sum is zero has probability at most
$2(1-m_\tau)^n\le2\exp(-nm_\tau)\to0$, because $\Prob(w_{a}>0)\ge\E w_a=m_\tau$.

\paragraph{Step 5 (variance estimator).} Write
$\xi_{ai}=w_{ai}(Y_i-\bar\mu_a)/m_\tau$. First, $n^{-1}\sum\xi_i^2/s_n^2$ has
mean one and variance at most $\E\xi^4/(ns_n^4)=O((n\tau)^{-1})$. Second,
\[
\hat\psi_{1i}-\xi_{1i}=w_{1i}(Y_i-\bar\mu_1)(1/\hat D_1-1/m_\tau)+w_{1i}(\bar\mu_1-\hat m_1)/\hat D_1 .
\]
The mean square of the first term is
$(n^{-1}\sum\xi_{1i}^2)(m_\tau/\hat D_1-1)^2=O_p(s_n^2)O_p((n\tau)^{-1})$.
For the second, $w_1^2\le w_1$ gives mean square at most
$(\hat m_1-\bar\mu_1)^2/\hat D_1$; by Step 4 and $\hat D_1=m_\tau(1+o_p(1))$
this is $O_p(1/(nm_\tau^2))$, which is $O_p((n\tau)^{-1})$ relative to
$s_n^2\asymp\tau^{-1}$. Arm $0$ is identical. By Minkowski's and the
Cauchy--Schwarz inequalities, $|n^{-1}\sum\hat\psi_i^2-n^{-1}\sum\xi_i^2|=o_p(s_n^2)$,
which gives (iii).

\paragraph{Step 6 (centering at $\beta_0$).}
$(\beta_{\ov,\tau}-\beta_0)/(s_n/\sqrt n)=O(\tau)O((n\tau)^{1/2})=O((n\tau^3)^{1/2})$.
With (i), (iii) and Slutsky's lemma this gives (iv).

\paragraph{Step 7 (exploration loss).} The non-preferred probability is
$q_\tau(h)=\sgm(-|h|/\tau)$. On $|h|\le u_0$,
$n\E[\phi(|H|)\sgm(-|H|/\tau)]\le n\bar\kappa\bar f\cdot2\tau^2\int_0^\infty
v\sgm(-v)dv=n\bar\kappa\bar f\tau^2\pi^2/6$, and the rest is at most
$n\bar G\exp(-u_0/\tau)$. In the strong form, on $v\le u_0/\tau$,
$\phi(\tau v)/\tau\to\kappa v$ and $f(\tau v)\to f(0)$ with dominating function
$\bar\kappa v\bar f\sgm(-v)$, which gives the asymptotic equality. For
$\tau_n=n^{-\zeta}$: $n\tau_n\to\infty$ iff $\zeta<1$; $n\tau_n^3\to0$ iff
$\zeta>1/3$; $n\tau_n^2\to0$ iff $\zeta>1/2$; and $n\exp(-u_0n^\zeta)\to0$.
The expected number of non-preferred actions is
$n\E[q_\tau(H)]=n\tau\int f(\tau u)\sgm(-|u|)du=2\log2\,f(0)\,n\tau(1+o(1))$ by the
same dominated convergence argument. \qed

\paragraph{Smoother effects.} If $|c(h)-c(0)-c'(0)h|\le M_2h^2$ and $f$ is
$L_f$-Lipschitz near $0$, then
$\E[\omega H]=\tau^2\int u\Omega(u)\{f(\tau u)-f(0)\}du=O(\tau^3)$ up to an
exponentially small tail, and the quadratic remainder contributes at most
$M_2\bar f\tau^3\int u^2\Omega/m_\tau$; the bias is $O(\tau^2)$. Differentiability of
$c$ alone is not enough: a bounded effect equal to $c(h)=1+|h|^{3/2}$
near zero gives bias
$\tau^{3/2}\int|u|^{3/2}\Omega(u)du\,(1+o(1))$.

\subsection{Proof of Proposition~\ref{prop:gapshape}}

(a) Substituting $h=\tau v$ on $|h|\le u_0$ and using $\sgm(-v)\le\exp(-v)$ elsewhere,
\[
R_n\le n\bar\kappa\bar f\int_{|h|\le u_0}|h|^p\sgm(-|h|/\tau)\,dh+n\bar G\exp(-u_0/\tau)
\le2\bar\kappa\bar fC_p\,n\tau^{p+1}+n\bar G\exp(-u_0/\tau).
\]
Here $C_p=\Gamma(p+1)(1-2^{-p})\zeta_R(p+1)$ with $\zeta_R$ the Riemann zeta function,
which follows from expanding $\sgm(-v)=\sum_{k\ge1}(-1)^{k+1}\exp(-kv)$; for example
$C_{1/2}\approx0.678$, $C_1=\pi^2/12$ and $C_2\approx1.803$. Validity of the interval
uses only (D1)--(D4) and $1/3<\zeta<1$, by Theorem~\ref{thm:boundary}(iv), and
$n\tau_n^{p+1}\to0$ iff $\zeta>1/(p+1)$.

(b) Restricting to $0<|H|\le u_0$ and substituting $h=\tau u$,
\[
R_n\ge n\lambda_0\tau\int_{|u|\le u_0/\tau}f(\tau u)\sgm(-|u|)\,du
=2\log2\,\lambda_0f(0)\,n\tau(1+o(1)),
\]
by dominated convergence with dominating function $\bar f\sgm(-|u|)$. \qed

%-----------------------------------------------------------------------------
\section{Proofs for Section~\ref{sec:population}}\label{app:population}

\paragraph{Proof of Theorem~\ref{thm:info}.}
\emph{Submodel.} For $\vartheta\in[-B,B]^J$ set $\mu^\vartheta_{a^*(h)}(h)=0$,
and on $S_j$ set $\mu^\vartheta_{\bar a(h)}(h)=s_j\vartheta_j$ with $s_j=+1$ if
$\bar a=1$ on $S_j$ and $s_j=-1$ otherwise; elsewhere
$\mu^\vartheta_{\bar a(h)}(h)=0$. Then $\mu^\vartheta\in\M(B,\sigma)$ and
$\theta(\mu^\vartheta)=\sum_jf_j\vartheta_j$.

\emph{Information.} The sample density is
$\prod_idF(H_i)\,\pi_i(A_i\mid H_i,\mathcal F_{i-1})\,\varphi_\sigma(Y_i-\mu^\vartheta_{A_i}(H_i))$,
where the assignment kernels $\pi_i$ are known functions of observed data and do
not depend on $\vartheta$. The score for $\vartheta_j$ is
$\sum_i\mathbf 1\{H_i\in S_j,A_i=\bar a(H_i)\}s_j(Y_i-s_j\vartheta_j)/\sigma^2$.
Its summands are martingale differences, so the Fisher information matrix is
diagonal with $I_j(\vartheta)=\E_\vartheta N_j/\sigma^2$, where
$N_j=\sum_i\mathbf 1\{H_i\in S_j,A_i=\bar a(H_i)\}$.

\emph{van Trees.} Take the product prior with marginal density
$B^{-1}\cos^2(\pi x/(2B))$ on $[-B,B]$, whose Fisher information is $J_0=\pi^2/B^2$.
With $\bar N_j=\int\E_\vartheta N_j\,\lambda(d\vartheta)$, the multivariate van
Trees inequality \cite{gill1995applications}, maximized over the direction
vector, gives
\[
\sup_{\M(B,\sigma)}\E(\hat\theta-\theta)^2\ge\sum_j\frac{f_j^2}{\bar N_j/\sigma^2+J_0}=\sum_j\frac{f_j\sigma^2}{a_j},
\qquad a_j=\frac{\bar N_j+\sigma^2J_0}{f_j}.
\]

\emph{Cauchy--Schwarz and loss.} Writing
$\sigma\sum_jf_j\sqrt{g_j}=\sum_j\sqrt{f_jg_ja_j}\cdot\sigma\sqrt{f_j/a_j}$,
\[
\sigma^2\Big(\sum_jf_j\sqrt{g_j}\Big)^2\le\Big(\sum_jf_jg_ja_j\Big)\Big(\sum_jf_j\sigma^2/a_j\Big).
\]
Here $\sum_jf_jg_ja_j=\sum_jg_j\bar N_j+\sigma^2J_0\sum_jg_j$, and since $g\ge g_j$
almost surely on $S_j$, $\sum_jg_jN_j\le\sum_ig(H_i)\mathbf 1\{A_i=\bar a(H_i)\}$
almost surely, so $\sum_jg_j\bar N_j\le\sup R_n$. This proves \eqref{eq:info}.

\emph{Consequences.} With $J=1$ and $S_1=\{g\ge\gamma\}\cap\{a^*=j\}$ of positive
mass for some $j$, \eqref{eq:info} gives
$\sup R_n\ge g_1\{\sigma^2f_1^2/\sup\E(\hat\theta-\theta)^2-\sigma^2\pi^2/B^2\}\to\infty$.
For the constant, fix a partition into level sets of $g$ intersected with the
two preferred-action sides; \eqref{eq:info} gives
$\liminf\delta_n^2\sup R_n\ge\sigma^2(\sum f_j\sqrt{g_j})^2$, and the lower sums
$\sum F(S)\essinf_S\sqrt g$ increase to $\E\sqrt g$ as the level mesh shrinks. \qed

\paragraph{Proof of Proposition~\ref{prop:achieve}.} The Horvitz--Thompson
estimator is unbiased with variance at most
$n^{-1}C\E[1/e+1/(1-e)]=n^{-1}C\E[1/q+1/(1-q)]\le n^{-1}C(\E[1/q]+2)$,
since $\{e,1-e\}=\{q,1-q\}$ and $q\le1/2$. Pointwise $1/q=\max\{2,\sqrt g/c_n\}\le\sqrt g/c_n+2$,
and the choice of $c_n$ gives $\E[1/q]+2\le n\delta_n^2/C$. The
loss is $n\E[gq]\le nc_n\E\sqrt g$. For uniform mixing with
$q=C/(n\delta_n^2-2C)\le1/2$, the same variance bound
gives $C(1/q+2)/n=\delta_n^2$, and the loss is $nq\E[g]=C\E[g]\,n/(n\delta_n^2-2C)$.
Without the correction, $q=C/(n\delta_n^2)$ gives risk $C/\{nq(1-q)\}>\delta_n^2$
at $\mu_0=\mu_1\equiv B$, where the bound is attained. \qed

%-----------------------------------------------------------------------------
\section{Proof of Theorem~\ref{thm:separation}}\label{app:separation}

The function $b$ satisfies $0\le b\le1$, is $\pi/(u_2-u_1)$-Lipschitz on
$\mathbb R$, vanishes outside $[u_1,u_2]$, and has $\E b(H)>0$ because $H$ has a
density and $F([u_1,u_2])>0$. Hence $|tb|\le B$ and $tb$ is $L$-Lipschitz for
$|t|\le t_{\max}$, so $\mu^t\in\Mlip$; $c^t=tb$, $\beta_0(\mu^t)=0$, and
$\theta(\mu^t)=t\E b(H)$. Every observation affected by the perturbation has
$A=0$ and $H\in[u_1,u_2]\subset(0,\infty)$, where $a^*=1$, so it is a
non-preferred action. Let $N=\sum_i\mathbf 1\{A_i=0,H_i\in[u_1,u_2]\}$.

\paragraph{(S1): uniform boundary coverage.} The loss statement is
Theorem~\ref{thm:boundary}(v), since the design does not depend on $\mu$. For
coverage we show that the error terms in the proof of Theorem~\ref{thm:boundary}
are bounded by constants that do not depend on $\mu\in\Mlip$. Here
$\sigma_a^2\equiv\sigma^2$, and $m_\tau$, $\hat D_a$ and the event of a zero weight
sum depend only on $(H_i,A_i)$, whose law does not depend on $\mu$. Throughout,
$C$ denotes a constant depending only on $(B,L,\sigma,\bar f)$.

\emph{Variance bounds.} Since $(1-e)+e=1$ and $|\mu_a-\bar\mu_a|\le2B$,
\eqref{eq:sn} gives
\[
\sigma^2/m_\tau\le s_n^2\le(\sigma^2+4B^2)/m_\tau .
\]

\emph{Berry--Esseen.} Given $(A,H)$, $Y-\bar\mu_a$ is Gaussian with variance
$\sigma^2$ and mean bounded by $2B$, so $\E[|Y-\bar\mu_a|^3\mid A,H]\le C$ and, as
$w_a^3\le w_a$ and $\E w_a=m_\tau$, $\E|\xi|^3\le Cm_\tau^{-2}$. The Berry--Esseen
inequality for independent and identically distributed summands gives
\[
\sup_x\Big|\Prob_\mu\Big(\frac{\sum_i\xi_i}{\sqrt ns_n}\le x\Big)-\Phi(x)\Big|
\le\frac{C\,\E|\xi|^3}{\sqrt n\,s_n^3}\le\frac{C}{\sigma^3\sqrt{nm_\tau}},
\]
uniformly in $\mu$.

\emph{Remainder.} In \eqref{eq:decomp}, $\Var(T_a)\le\E[w_a^2(Y-\bar\mu_a)^2]/(nm_\tau^2)
\le(\sigma^2+4B^2)/(nm_\tau)$, and $\Var(\hat D_a/m_\tau)\le1/(nm_\tau)$ does not
depend on $\mu$. By Chebyshev's inequality, for every $\varepsilon>0$ there is $M$
with $\sup_\mu\Prob_\mu(|T_a|>M(nm_\tau)^{-1/2})<\varepsilon$ and
$\Prob(|m_\tau/\hat D_a-1|>M(nm_\tau)^{-1/2})<\varepsilon$ for large $n$. Dividing by
$s_n/\sqrt n\ge\sigma(nm_\tau)^{-1/2}$, the remainder in units of $s_n/\sqrt n$ is
$O_p((nm_\tau)^{-1/2})$ uniformly in $\mu$.

\emph{Variance estimator.} Each bound in Step 5 of the proof of
Theorem~\ref{thm:boundary} involves only $\E\xi^4\le Cm_\tau^{-3}$, $s_n^2\ge\sigma^2/m_\tau$,
the bounds on $T_a$ above and the law of $\hat D_a$; hence
$\sup_\mu\Prob_\mu(|\hat s^2/s_n^2-1|>\varepsilon)\to0$ for every $\varepsilon>0$.

\emph{Bias.} Global Lipschitz continuity gives, with no tail term,
$|\beta_{\ov,\tau}-\beta_0|\le2L\E[\omega|H|]/m_\tau\le4\log2\,L\bar f\tau^2/m_\tau$,
so in units of $s_n/\sqrt n$ the bias is at most
$C\tau^2\sqrt{n/m_\tau}=O((n\tau^3)^{1/2})$ uniformly in $\mu$, using
$m_\tau=\tau f(0)(1+o(1))$.

\emph{Coverage.} Write $Z_n=(\hat\beta-\beta_{\ov,\tau})\sqrt n/s_n$ and
$b_n=(\beta_{\ov,\tau}-\beta_0)\sqrt n/s_n$. The interval covers $\beta_0$ iff
$|Z_n+b_n|\le z_{1-\eta/2}\hat s/s_n$. By the three displays above, for every
$\varepsilon>0$ there are events of $\Prob_\mu$-probability at least $1-\varepsilon$
for all $\mu$ and large $n$, on which $|Z_n-n^{-1/2}\sum_i\xi_i/s_n|\le\varepsilon$,
$|b_n|\le\varepsilon$ and $|\hat s/s_n-1|\le\varepsilon$. Hence the coverage
probability lies within $\varepsilon+C/\sqrt{nm_\tau}+\sup_{|x|\le z+3\varepsilon}\varphi(x)\cdot6\varepsilon$
of $1-\eta$ for all $\mu\in\Mlip$, where $\varphi$ is the standard normal density,
and letting $\varepsilon\downarrow0$ proves (S1).

\paragraph{(S2).} The sample densities under $\mu^0$ and $\mu^t$ coincide on
$E=\{N=0\}$, so $P_0(E)=P_t(E)$, and for any event $A$,
$|P_0(A)-P_t(A)|=|P_0(A\cap E^c)-P_t(A\cap E^c)|\le\max\{P_0(E^c),P_t(E^c)\}=P_0(E^c)$.
Thus $\TV\le P_0(N\ge1)\le\E_0N$. Because $g\ge\gamma$ $F$-almost surely on
$[u_1,u_2]$ and $H_i\sim F$, $\gamma N\le\sum_ig(H_i)\mathbf 1\{A_i=\bar a(H_i)\}$
almost surely, so $\E_0N\le R_n(\mu^0)/\gamma$. Le Cam's two-point argument with
separation $t_1\E b$ gives the displayed bound.

\paragraph{(S3).} The score of $t\mapsto\mu^t$ is
$\sum_i\mathbf 1\{A_i=0\}b(H_i)\{-(Y_i+tb(H_i))\}/\sigma^2$, with information
$I(t)=\E_t\sum_i\mathbf 1\{A_i=0\}b(H_i)^2/\sigma^2\le\E_tN/\sigma^2\le R_n(\mu^t)/(\gamma\sigma^2)$.
With the prior $t_{\max}^{-1}\cos^2(\pi t/(2t_{\max}))$ on $[-t_{\max},t_{\max}]$,
whose information is $\pi^2/t_{\max}^2$, the van Trees inequality for
$\psi(t)=t\E b$ gives the bound. \qed

%-----------------------------------------------------------------------------
\section{Proofs and a lemma for Section~\ref{sec:temperature}}\label{app:temperature}

\paragraph{Proof of Theorem~\ref{thm:common}.} Take the single cell
$S=\{|\Delta|\ge D_{\max}-\epsilon,a^*=j\}$, with $j$ attaining $\nu_\epsilon$.
There $q_\tau\le \exp(-(D_{\max}-\epsilon)/\tau)$, so
$\bar N\le n\nu_\epsilon\exp(-(D_{\max}-\epsilon)/\tau)$, and the van Trees step
of Theorem~\ref{thm:info} with $J=1$ gives
$\delta_n^2\ge \nu_\epsilon^2\sigma^2/(n\nu_\epsilon\exp(-(D_{\max}-\epsilon)/\tau)+\sigma^2\pi^2/B^2)$.
Once $\delta_n^2\le \nu_\epsilon^2B^2/(2\pi^2)$, this forces
$n\nu_\epsilon\exp(-(D_{\max}-\epsilon)/\tau)\ge \nu_\epsilon^2\sigma^2/(2\delta_n^2)$, that is,
$\tau_n\ge(D_{\max}-\epsilon)/\log\{2n\delta_n^2/(\nu_\epsilon\sigma^2)\}$, where the
logarithm is positive because $2n\delta_n^2>\nu_\epsilon\sigma^2$. Nothing here
requires $\delta_n\to0$. On the other hand
\[
R_n=n\E[\phi(|\Delta|)\sgm(-|\Delta|/\tau)]
\ge n\underline\kappa f_{\min}\tau^2\int_0^{v_0/\tau}v\sgm(-v)\,dv
\ge n\underline\kappa f_{\min}c_*\min\{\tau,v_0\}^2 .
\]
Combining the two bounds proves the claim. With finitely many score values, the
same argument gives $\tau\ge\Delta_{\max}/\log(C'n\delta^2)$ for $\delta$ small
enough, and the $\Delta_{\min}$ cell, with mass $f$, contributes at least
$nf\phi(\Delta_{\min})\exp(-\Delta_{\min}/\tau)/2$, which is of order
$n^{1-\Delta_{\min}/\Delta_{\max}}$ up to logarithmic factors. \qed

\paragraph{Remark~\ref{rem:mixture}.} Write $u=|H|$, $C=\sigma^2+B^2$,
$r(u)=\sgm(-u)$, and $q_n=w_nr+(1-w_n)\sgm(-nu)\le1/2$ for the mixture. Since
$q_n\ge w_nr$, $\E[1/q_n]\le\E[1/r]/w_n=n\delta^2/C-2$, and the variance bound in
the proof of Proposition~\ref{prop:achieve} gives uniform risk at most
$C(\E[1/q_n]+2)/n\le\delta^2$. The loss is
$n\E[uq_n]\le nw_n\E[ur]+n\E[u\sgm(-nu)]=C\E[1/r]\E[ur]\,n/(n\delta^2-2C)+O(1/n)$,
which equals $C\E[1/r]\E[ur]/\delta^2+O(1/n)$.

\paragraph{Proof of Proposition~\ref{prop:ctx}.} Solving
$\sgm(-|\Delta|/\tau)=q$ gives the formula. If $g(h)=0$, then $q_n(h)=1/2$, which
is implementable by $\tau=\infty$ (and by any $\tau$ when $\Delta(h)=0$); if
$\Delta(h)\ne0$ and the cap binds, $q_n(h)=1/2$ again requires $\tau=\infty$. If
$\Delta(h)\neq0$ and $g(h)>0$, then $q_n(h)=c_n/\sqrt{g(h)}<1/2$ once
$c_n<\sqrt{g(h)}/2$, and
$\log\{(1-q_n)/q_n\}=\log(1/c_n)+\tfrac12\log g(h)+\log(1-q_n)$ with
$c_n\to0$, which gives the limit. \qed

\paragraph{Misspecified gaps.} Suppose the allocation uses $\hat g$ in place of $g$.
Designs are compared at a common value of the sparse-exploration criterion
$\sigma^2\E[1/q]/n$, which is an allocation criterion, not the mean squared error
of an estimator. Let $\Lambda(\rho)=(1+\sqrt\rho)^2/(4\sqrt\rho)$.

\begin{lemma}\label{lem:robust}
Assume $0<\E\sqrt g<\infty$ and $0<a\le\hat g/g\le b<\infty$ on $\{g>0\}$, with
$\hat g=0$ on $\{g=0\}$, and let $\rho=b/a$.
\begin{enumerate}[label=(\alph*),leftmargin=2em,itemsep=1pt]
\item If $F(g>0)=1$ and the uncapped designs $q=c/\sqrt{\hat g}$ and $q=c'/\sqrt g$,
calibrated to the same criterion, are at most $1/2$ almost surely, the ratio of
their losses is at most $\Lambda(\rho)$. The bound is attained when the measure
$w\propto\sqrt g\,dF$ has a set of mass exactly one half.
\item For the capped designs $q_c=\min\{1/2,c/\sqrt{\hat g}\}$ and
$q'_{c'}=\min\{1/2,c'/\sqrt g\}$, equal to $1/2$ where the gap is zero and
calibrated to a common value $T=\E[1/q_c]=\E[1/q'_{c'}]$, the loss ratio converges,
as $T\to\infty$, to $\E[\sqrt{\hat g}]\E[g/\sqrt{\hat g}]/(\E\sqrt g)^2\le\Lambda(\rho)$.
\end{enumerate}
\end{lemma}

If $\essinf g=0$, the uncapped design exceeds $1/2$ on a set of positive mass for
every $c>0$, so part (a) is vacuous and part (b) is the relevant statement; this
includes $g(h)=|h|$. A common rescaling of $\hat g$ changes nothing. Gap-based
allocation beats uniform mixing for every admissible $\hat g$ if
$\E[g]/(\E\sqrt g)^2>\Lambda(\rho)$.

\paragraph{Proof of Lemma~\ref{lem:robust}.} (a) With $q=c/\sqrt{\hat g}$ and
$\E[1/q]=T$, $c=\E\sqrt{\hat g}/T$ and the loss is
$n\E[\sqrt{\hat g}]\E[g/\sqrt{\hat g}]/T$, against $n(\E\sqrt g)^2/T$ for $\hat g=g$.
With $s=(\hat g/g)^{1/2}\in[\sqrt a,\sqrt b]$ and the probability measure
$w\propto\sqrt g\,dF$, the ratio is $\E_w[s]\E_w[1/s]$, which Kantorovich's
inequality bounds by $(\sqrt a+\sqrt b)^2/(4\sqrt{ab})=\Lambda(b/a)$. If $w$ has a set
of mass one half, setting $s=\sqrt a$ there and $s=\sqrt b$ elsewhere attains the
bound. With atoms the bound need not be attained: for $F=(0.99,0.01)$, $g=(1,100)$
and $\rho=36$, uniform mixing costs $1.675$ times the optimum, below
$\Lambda(36)=2.042$, yet every admissible $\hat g$ costs at most $1.347$ times the
optimum.

(b) By dominated convergence, $c\,\E[1/q_c]=\E\max\{2c,\sqrt{\hat g}\}\to\E\sqrt{\hat g}$
and $\E[gq_c]/c\to\E[g/\sqrt{\hat g}]$, with $g/\sqrt{\hat g}=0$ where $g=0$, as
$c\to0$; likewise for $q'$. Calibration to a common $T\to\infty$ forces $c,c'\to0$
with $cT\to\E\sqrt{\hat g}$ and $c'T\to\E\sqrt g$, so the loss ratio
$\E[gq_c]/\E[gq'_{c'}]$ converges to the stated limit, which is bounded as in (a). \qed
